% 本文件由 LaTeX 排版 Agent 根据有效 translation.json 自动生成。
% author=Miscellaneous; work_id=0858; title=沙穆纳、古里亚与哈比卜殉道录

\chapter{\Person{沙穆纳}、\Person{古里亚}与\Person{哈比卜}殉道录}

% structure: p0001 source=h1 target=omitted reason=page-title-already-rendered-by-parent

\Emph{圣宣信者\Person{沙穆纳}、\Person{古里亚}与\Person{哈比卜}殉道录，取自\Person{西默盎·默塔弗拉斯特斯}。}\par

自马其顿人亚历山大帝国算起第六百年，当\Person{戴克里先}为罗马君主第九年，\Person{马克西米安}第六次任执政官，\Person{奥加尔}之子\Person{奥加尔}任总督，\Person{科格纳图斯}为厄德萨主教时，一场大迫害兴起，波及罗马治下诸国之教会。基督徒之名被视为可憎，遭受攻击与诽谤；司铎与隐修士因坚贞不屈，备受骇人刑罚，虔敬者因忧愁恐惧而不知所措。他们因对基督的挚爱渴望宣扬真理，却又因畏惧刑罚而退缩。因为那些武装起来反对真宗教者，一心要迫使基督徒背弃基督信仰，皈依农神与\Person{瑞亚}之教；而信友则努力证明异教崇拜的对象实不存在。\par

此时，有人向法官控告\Person{古里亚}与\Person{沙穆纳}。前者生于\Place{萨尔西吉图阿}，后者生于\Place{加纳斯}村；然二人均在\Place{厄德萨}长大——此地他们称为\Place{美索不达米亚}，因其位于\Place{幼发拉底河}与\Place{底格里斯河}之间。此城此前鲜为人知，但经其殉道者之奋斗后，名扬天下。这些圣者决不愿居于城中，反而远离尘嚣，如同愿远离纷扰者，只求显扬于天主面前。\Person{古里亚}的纯洁与仁爱是他珍贵而可敬的财富；由于他修养前者，得号\Quote{纯洁者}，以致你从名字不知其为谁，惟以绰号呼之。\Person{沙穆纳}将自己的身体、青春活泼的心灵奉献于天主，在德性上堪与\Person{古里亚}媲美。有人向法官控告他们，说他们不仅以其教诲遍及\Place{厄德萨}四周，鼓励民众持守信仰，而且引导他们藐视迫害者，并教导他们完全蔑视不虔敬之事，正如经上所记：\BibleQuote{不要依赖王侯——依赖人子，其中并无救援。}这些陈述使法官狂怒至极，下令将所有尊崇基督宗教、追随\Person{沙穆纳}与\Person{古里亚}之教训者，连同那些劝诱他们如此行的人，一并逮捕，严密关押。命令得以执行；他便趁机鞭挞其中一些人，用各种方式折磨另一些人，迫使他们服从皇帝之命；然后，仿佛仁慈宽厚，他允许其他人回家；但\Emph{我们这两位}圣人，作为首倡者并将虔敬传与他人者，他下令在狱中继续虐待。然而他们因殉道的共融而喜乐。因为他们听说其他省份有许多人也经历了同样的战斗：在\Place{凯撒勒雅}的\Place{巴勒斯坦}，有\Person{埃皮法尼乌斯}、\Person{伯多禄}与至圣的\Person{庞菲鲁斯}，及许多其他人；在\Place{迦萨}，\Person{弟茂德}；在\Place{亚历山大}，大\Person{弟茂德}；在\Place{得撒洛尼}，\Person{阿加佩图斯}；在\Place{尼科美底亚}，\Person{赫西基乌斯}；在\Place{阿德里安堡}，\Person{斐理伯}；在\Place{梅利蒂纳}，\Person{伯多禄}；\Person{赫尔墨斯}及其同伴在\Place{殉道者城}边界：所有这些人都由公爵\Person{赫拉克利安努斯}戴上了殉道者的冠冕，还有无数其他宣信者，我们无法尽知。但我们必须回到先前所说之事。\par

于是，\Place{厄德萨}总督\Person{安多尼}允许其他人回家后，设立了一个高高的审判席，命令将殉道者带到他面前。侍从遵命而行，总督对圣人们说：\Quote{我们最神圣的皇帝命令你们放弃所持守的基督信仰，并在祭台上向\Person{朱庇特}焚香，以行神明之礼。}\Person{沙穆纳}回答说：\Quote{我们决不可背弃真信仰，因我们借此希望获得不朽，反而去崇拜人手所造之物和偶像！}总督说：\Quote{皇帝的命令必须绝对服从。}\Person{古里亚}答道：\Quote{我们绝不否认我们纯洁神圣的信仰，去追随终将朽坏之人的意志。因为我们有一位天上的父，我们遵行祂的旨意，祂说：\BibleQuote{凡在人前承认我的，我在天上的父前也必承认他；但谁在人前否认我的，我在我的父及祂的天使前也必否认他。}}审判官说：\Quote{那么你们拒绝服从皇帝的旨意？你们岂能以为，普通人的意图——那些与你们同样无权势的人——真能实现，而拥有最高权力者的命令却落空？}圣人们说：\Quote{凡遵行万王之王旨意的人，就厌弃并拒绝肉身的旨意。}总督以死威胁他们，除非他们服从，\Person{沙穆纳}说：\Quote{暴君，我们若遵行造物主的旨意，就必不死；不但不死，反而活着；但若我们服从你皇帝的命令，须知即使你\Emph{不}处死我们，我们同样会悲惨地灭亡。}\par

总督闻言，即命狱吏\Person{阿诺维图斯}将二人严加看管。因为本性倾向邪恶的心灵无法忍受真理，正如病眼无法忍受太阳的明亮光芒。狱吏遵命而行，殉道者被囚于狱中，那里先前已有许多其他圣人被士兵关押。此时，皇帝\Person{戴克里先}召来\Place{安提约基雅}总督\Person{穆索尼乌斯}，命他前往\Place{厄德萨}，查看被囚禁在那里的基督徒，无论他们是平民还是圣职人员，都要审问他们的信仰，并按他认为合适的方式处置他们。于是穆索尼乌斯来到\Place{厄德萨}，将\Person{沙穆纳}和\Person{古里亚}首先置于审判台前，对他们说：\Quote{天下之主谕令，你们必须献上酒奠，并在朱庇特祭台上献香。若你们拒绝，我将用多种刑罚毁灭你们：我将用鞭子将你们的身体撕裂，直至你们的脏腑；我将不断将滚烫的铅液倒入你们的腋窝，直至流到你们的肠腹；之后，我将把你们悬挂起来，时而吊手，时而吊脚，并放松你们关节的绑索；我还要发明前所未知的新刑罚，叫你们完全无法忍受。}\par

\Person{沙穆纳}答道：\Quote{我们畏惧\BibleQuote{那虫}，祂的威胁是针对否认主的人而宣告的，也畏惧\BibleQuote{那永不熄灭的火}，甚于你摆在我们面前的这些酷刑。因为\Emph{天主}自己——我们所献上合理事奉的对象——首先会赐我们力量忍受这许多酷刑，并将我们从你手中拯救出来；之后，还要使我们安息于安全之所，那里是所有喜乐者的居处。此外，你所攻击的不过是身体而已：你能对灵魂造成什么伤害呢？因为灵魂住在身体内时，它凌驾于酷刑之上；而当它离开时，身体便毫无感觉。因为\BibleQuote{我们外在的人虽被摧毁，内在的人却日日更新}，正是藉着忍耐，我们通过这场摆在我们面前的竞赛。}然而，总督又以一种宣告的姿态，为了使他们若不顾从，便更能公正地惩罚他们，说道：\Quote{请放弃你们的错误吧，我恳求你们，服从皇帝的命令吧：你们将无法忍受这些酷刑。}圣\Person{古里亚}答道：\Quote{我们并非如你所说，是错误的奴隶；我们永远不会服从皇帝的命令。天主禁止我们如此懦弱无知！因为我们是祂的门徒，祂为我们舍弃了生命，以此彰显祂的良善与慈爱的丰盛。因此，我们甚至抵抗罪恶至死；无论如何，我们都不会被仇敌的诡计击败——第一个人曾因此被诱骗，因悖逆而从树上摘取了死亡；\Person{加音}曾被说服，染手足兄弟之血后，发现罪恶的酬报只是哀号和恐惧。但我们听从基督的话，将\BibleQuote{不要害怕那些杀害身体却不能杀害灵魂的}，反而要害怕\BibleQuote{那能毁灭我们灵魂和身体的那位}。}暴君说：\Quote{我压抑愤怒、显忍耐，并非为了让你们有机会引用你们自己的著作来反驳我的指控，而是为了让你们执行皇帝的命令，平安返回家园。}\par

这些话丝毫未能动摇殉道者的决心；但他们走近一步，说：\Quote{这于我们何干？即使你\Emph{是}发怒，鼓怒，如雪花般将酷刑倾降在我们身上，那岂不更是厚待我们？因为那样会使我们坚韧的证明更加显著，也为我们赢得更大的赏报。因为这是我们希望的高峰：我们将离开这暂时的居所，往永远的居所去。因为我们在天上有一座\BibleQuote{非人手所造的帐幕}，圣经通常也称它为\BibleQuote{\Person{亚巴郎}的怀抱}，因为亚巴郎蒙受与天主亲密交往之福。}总督见他们毫不动摇，立即停止言语，施行所威胁的刑罚，命狱吏\Person{阿努伊努斯}将二人分别用一只手悬挂起来；当他们的手因承受全身重量而脱臼时，再在他们脚上挂一块重石，使痛苦更加剧烈。这命令被执行了，从第三时辰到第八时辰，他们刚毅地忍受着这剧烈酷刑，不发一言，不哼一声，也没有任何懦弱或卑贱的表现。你会以为他们是在不属于自己的身体上受苦，或是别人在受苦，而他们不过是旁观者。\par

同时，当他们被吊着手悬挂时，总督正在审理其他案件。处理完这些案件后，他命令狱卒询问这些圣人是否愿意服从皇帝并从酷刑中获释。当狱卒向他们提出这个问题时，发现他们要么不能要么不愿回答，他就命令把他们关在监狱内室一个黑暗的地牢里——名副其实的黑暗，并把他们的脚锁在足枷中。黎明时分，他们的脚从足枷的束缚中松开；但监狱被严严关闭，连一丝阳光都透不进来；狱卒奉命在整整三天内不给他们一点面包或一滴水。因此，除了其他一切之外，殉道者还被判处在黑暗的监狱中并长期禁食。到了第三天，大约八月初，监狱被打开让光线进入，但他们仍被关在里面直到十一月十日。然后审判官把他们带到他面前，说道：\Quote{难道这所有的时间还不足以使你们改变主意，作出某种有益的决定吗？}他们回答说：\Quote{我们已经多次告诉过您我们的心意；因此，请执行您所受的命令。}总督立即命令\Person{沙姆纳}单膝跪下，并在他的膝盖上系上一条铁链。做完这事后，他把他头朝下吊起来，用让他跪下那只脚悬挂，另一只脚则用一块无法言喻的重铁向下拉，企图把这位斗士撕裂成两半。因此，髋骨窝被扯出原位，\Person{沙姆纳}成了瘸子。然而，对于\Person{古里亚}，因为他体弱且面色苍白，他未加惩罚——不是因为他以友善的目光看待他，也不是因为同情他的软弱，而是为了把他留到另一个机会再来惩罚这个他急于惩罚的人，免得正像我说的那样，由于我的疏忽，他在承受为他预备的酷刑之前就耗尽了。\par

这时，自\Person{沙姆纳}被吊起以来，白天已过了两个时辰，第五个时辰也已到来，他仍然被悬挂在高处——周围的士兵起了怜悯之心，劝他服从皇帝的命令。但罪人的同情对这位圣人毫无影响。因为虽然他因酷刑而痛苦不堪，但他没有给他们任何回答，任由他们随意哀悼，自认为他们自己而不是他才值得怜悯。然而，他举目望天，从内心深处向天主祈祷，提醒祂古时所行的奇事：他说道：\Quote{上主天主，没有祢，连一只可怜的小麻雀也不会落入网罗；是祢在痛苦中安慰了\Person{达味}的心；是祢赐予\Person{达尼尔}力量以对抗狮子；是祢使\Person{亚巴郎}的子孙战胜了暴君和火焰。现在，上主，求祢也垂顾这场针对我们的战争，因为祢熟知我们本性的软弱。因为仇敌企图将祢右手的化工从与祢同在的荣耀中夺走。但求祢以慈悲的目光注视我们，在我们内保持祢诫命的明灯，抵挡一切试图熄灭它的企图；并以祢的光引导我们的道路，赐予我们享受在祢内的幸福，因为祢是永受赞颂的，万世无疆。}就这样，他赞颂了竞赛的裁判者；在场的一位书记员将他的话记录下来。\par

最后，总督命令狱卒解除对他的惩罚。狱卒照做了，把他带走，他因所受的痛苦而昏厥无力，他们把他送回原来的监狱，让他躺在\Person{古里亚}旁边。然而，十一月十五日夜间，大约鸡叫的时候，审判官起身。火炬和随从走在他前面；到达所谓的巴西利卡（即法庭所在）时，他隆重地登上审判席，派人去请两位斗士\Person{古里亚}和\Person{沙姆纳}。后者在两个\Emph{狱卒们的}搀扶下走进来，由两人搀扶：因为他饥饿疲惫，年迈体弱，只有美好的希望支撑着他。\Person{古里亚}也必须被抬进去：因为他完全不能行走，因为他的脚被链子严重磨伤。不虔敬的辩护人对他们两人说：\Quote{根据所给予的许可，你们\Emph{无疑}已经共同商议过你们该做之事。那么告诉我，你们是否达成了任何新的决定？你们是否在任何方面改变了先前的心意？服从最神圣的\Emph{皇帝}的命令吧！因为这样你们将恢复享受你们的财产和所有物，甚至这最令人愉快的光明。}对此，殉道者们回答说：\Quote{没有一个明智的人会看重暂时享受那些转瞬即逝的东西。我们过去使用和看见它们的时间已经足够；我们并不感到缺少任何一样。相反，你们威胁我们的死亡将把我们带入不朽的居所，并使我们分享那边的幸福。}\par

\Quote{总督回答说：你所讲的话使我耳中充满极大的忧愁。不过，我要向你说明所决定的：如果你将乳香放在祭台上，向\OriginalTerm{朱庇特}{Jupiter}的像献祭，一切都会顺利，你们各自可以回家；但是，如果你们仍然固执违抗皇帝的命令，你们必定会失去头颅：因为伟大的皇帝正是这样意愿和决定的。} 对此，最高尚的\Person{沙穆纳}回答说：\Quote{如果你赐予我们如此大的恩惠，使我们脱离今生的痛苦，并放我们前往来世的福乐，那么，就我们所能，你必会从那位按我们的益处支配我们所有物的主宰那里得到赏报。} 总督似乎有点和蔼地回答说：\Quote{我至今一直忍耐，容忍了你那些长篇大论，为的是通过拖延，你能改变心意，走向对你有益的事，而不必遭受死亡的刑罚。} 他说：\Quote{那些为基督的缘故而甘愿承受暂时死亡的人，显然将脱离永死。因为死于世俗的人在基督内活着。因为\Person{伯多禄}，在宗徒团体中如此闪耀，被判处十字架和死亡；\Person{雅各伯}，雷霆之子，被\Person{黑落德·阿格黎帕}用刀杀死。此外，\Person{斯德望}也被石头砸死，他是第一位跑完殉道赛程的人。至于\Person{若翰洗者}，你又要说什么呢？你必定承认他那卓越的刚毅和直言不讳，他宁愿死也不对婚姻不忠保持沉默，而那个奸妇以她的舞蹈得到了他的头作为奖赏。}\par

\Quote{我如此忍耐你们，并不是为了让你们数算你们所谓的圣人，而是要你们改变决定，服从皇帝的命令，从而免于极痛苦的死亡。因为，如果你们表现得如此过分大胆和傲慢，你们还能指望什么，除了更严厉的刑罚在等着你们，在它们的压力下，你们即使不情愿也会做我要求你们做的事；到那时，求助于怜悯已为时过晚。因为被迫发出的呼号无权引起怜悯；而自愿发出的呼号才值得同情。} \Person{基督}的宣信者和殉道者说：\Quote{不需要多话。看！我们已准备好承受你加给我们的一切刑罚。因此，命令你做的事，不要迟延。因为我们是\Person{基督}真天主的崇拜者，（我们再说一遍）祂的国度没有终结；祂也唯独能酬报那些光荣祂名的人。} 与此同时，当圣人们说这些话时，总督判决他们应受刀剑之死。但他们充满了言语无法表达的喜乐，喊道：\Quote{光荣和赞美归于你，你是万有之神，因为你乐意让我们所进入的战斗进行到底，并且也让我们从你手中接受永不褪色的光明。}\par

当总督看到他们不屈不挠的坚定，以及他们如何以灵魂的欢欣听到了最后的判决时，他对圣人们说：\Quote{愿天主查考所行的，\Emph{并作见证}，就我而言，我决不愿意你们失去性命；但皇帝给我不可更改的命令迫使我这样做。} 然后他命令一名戟兵负责看管殉道者，将他们放在一辆车上，带一些士兵运到城外，在那里用刀结束他们。于是戟兵在深夜，当市民沉睡时，带着圣人们从\Place{罗马门}出去，将他们带到城北的\Place{贝特拉比克拉山}。到达那里后，他们怀着心中的喜乐和极大的坚定下了车，请求戟兵和那些听他指挥的人给他们时间祈祷；请求得到了准许。因为，仿佛他们的折磨和流血不足以替他们祈求，他们仍因自己的谦卑认为有必要祈祷。于是他们举目向天，恳切祈祷，以这句话作为结束：\Quote{天主和我们的主耶稣基督的父，求你在平安中接纳我们的灵魂到你那里。} 然后\Person{沙穆纳}转向戟兵说：\Quote{执行你所受的命令。} 于是他同\Person{古里亚}一起跪下，被斩首，时在11月15日。这就是关于殉道者所发生的事的记述。\par

但因所寻找的数目需要一个第三位，好使天主圣三在他们身上受光荣，奇妙的天主眷顾找到了\Person{哈比布}——虽然后来才寻得，但他也与先前那些先他而往的人一样，决定踏上旅程，并在他们殉道的同一日达到了圆满。\Person{哈比布}，这位伟大的殉道者，与他们是同乡，即\Place{特瑟亚}村人；且蒙受神圣执事职的尊荣。但当\Person{李锡尼}执掌罗马帝国的权柄，\Person{吕撒尼雅}被任命为\Place{厄德撒}总督时，针对基督徒的迫害再次兴起，普遍的危险威胁临到\Person{哈比布}。因为他走遍全城，向所遇见的每个人教导圣经，并勇敢地试图坚强他们于虔敬中。此事传到\Person{吕撒尼雅}耳中，他便向皇帝\Person{李锡尼}告发。因他渴望自己受托办理审讯基督徒的事，尤其是\Person{哈比布}：因他以前从未受托此事。于是皇帝派人送信给他，命令他处死\Person{哈比布}。\Person{吕撒尼雅}接信后，到处搜查\Person{哈比布}；\Person{哈比布}因在教会中的职务，住在城内某处，他的母亲和一些亲属与他同住。他得知此事后，唯恐因退出殉道的行列而受罚，便自动前往，向一位御林军首领名叫\Person{特奥特库诺斯}的人自首，随即说：\Quote{我就是你们所搜寻的\Person{哈比布}。} 但那人善意地看着他说：\Quote{好人，还没有人知道你来我处：所以走吧，保全你的安全；不必挂念你的母亲，也不必挂念你的亲属：因为他们不可能有任何麻烦。} \Person{特奥特库诺斯}这样说。\par

但\Person{哈比布}，因此时机正需要殉道，便不肯顺从软弱懦怯的心，以任何隐秘方式保全己命。因此他回答说：\Quote{我告发自己，并非为了亲爱的母亲，也非为了亲属；我来是为了宣认\Person{基督}。因为看！不论你同意与否，我都要出现在总督面前，并在王侯面前宣告我的\Emph{主}\Person{基督}。} \Person{特奥特库诺斯}因此担心他可能自动前往总督处，这样他自己就会因未告发他而陷入危险，于是抓住\Person{哈比布}，带他到总督面前说：\Quote{这就是被搜寻的\Person{哈比布}。} \Person{吕撒尼雅}得知\Person{哈比布}自动前来投入争战，便断定这是藐视和狂妄无礼的表示，仿佛他轻看审判席的威严；于是立即将他提审。他询问他的生活状况、姓名和籍贯。\Person{哈比布}答说自己原是\Place{特瑟亚}村人，并表明自己是\Person{基督}的仆役，总督随即斥责殉道者不服从皇帝的命令。他强调此事的明证是他拒绝向朱庇特献香。\Person{哈比布}一再回答说他是基督徒，不能离弃真\Person{天主}，也不能向\Emph{人手}所造、没有知觉的无生命之物献祭。总督于是下令用绳索绑住他的双臂，将他高高吊在梁上，并用铁爪撕裂他。悬吊之苦远比撕裂更难忍受：因他的双臂被强力拉伸，有被扯裂的危险。\par

在此期间，当他悬挂在空中时，总督转而使用温和的言词，并装出耐心的样子。但他继续以更严厉的刑罚威胁他，除非他改变决心。但他说：没有人能使我放弃信德，也没有人能说服我崇拜魔鬼，即使他施加更多更大的酷刑。总督问他，他期望从摧毁他全身的酷刑中得到什么好处时，基督的殉道者\Person{哈比卜}回答说：我们所关注的并非只限于眼前，我们也不追求看得见的事物；如果你也愿意将目光转向我们的望德和所应许的报酬，或许你甚至会同\Person{保禄}一起说：\BibleQuote{现时的苦难，与将来要在我们身上显示的光荣，是不能比较的。}总督断言他的话是愚蠢之言；他看到尽管他使出浑身解数，时而甜言蜜语并装出耐心，时而又以可怕的死亡威胁恫吓他，却无法以任何方式说服他时，他在宣判时说：我不会让你突然迅速死去；我将用慢慢燃烧的火，逐渐使你解体，从而让你收起那凶猛倔强的脾气。于是，在城外北边的一个地方，堆起了一些木柴，他被带到火堆旁，他的母亲以及与他有血缘关系的亲属跟随着他。他祈祷，祝福众人，并在主内与他们亲吻告别；之后，木柴被他们点燃，他被投入火中；当他张口接受火焰时，他将灵魂交还给了赐予它的一位。火势平息后，他的亲属用一块贵重的亚麻布包裹他，并用香膏涂抹他；他们合宜地咏唱圣咏和圣歌后，将他安放在\Person{沙慕纳}和\Person{古里亚}旁边，以光荣圣父、圣子及圣神，他们是不可分开的天主圣三：尊威与崇拜归于祂，现今，永远，直至无穷之世。阿们。这就是殉道者\Person{哈比卜}在\Person{里基纽}时代生命的终结，他由此获得了与圣人同葬的特恩，并为虔诚者带来了迫害中的安息。因为不久之后，\Person{里基纽}的势力衰落，\Person{君士坦丁}的统治兴盛，\Place{罗马}帝国的主权归于他；他是第一位公开表白虔诚的皇帝，允许基督徒按基督徒的方式生活。\par

\SourceRecord{Translated by B.P. Pratten.；Ante-Nicene Fathers；Vol. 8.；Edited by Alexander Roberts, James Donaldson, and A. Cleveland Coxe.；Buffalo, NY: Christian Literature Publishing Co.,；1886.；<http://www.newadvent.org/fathers/0858.htm>.}
